%══════════════════════════════════════════════════════════════════════════════
\section{Phase 4: Model Design \& Implementation}
%══════════════════════════════════════════════════════════════════════════════

\subsection{Architecture Overview}

The system consists of four modality-specific encoders whose outputs are fused by a learned gating mechanism. Two task heads produce the final predictions.

\begin{figure}[H]
\centering
\includegraphics[width=\textwidth]{figures/architecture.png}
\caption{Complete multimodal fusion architecture. The HAR-RV skip connection (red arrow) bypasses the CNN-BiLSTM bottleneck to directly feed autoregressive features into the fusion trunk.}
\end{figure}

\subsection{Model 1: CNN-BiLSTM Price Encoder}

\subsubsection{Design Rationale}

\begin{itemize}[noitemsep]
    \item \textbf{1D-CNN:} Detects local patterns in 3-day windows (e.g., momentum, reversals). Uses 64 filters to learn 64 different pattern types simultaneously.
    \item \textbf{BiLSTM:} Captures how CNN patterns evolve over the full 60-day window. Bidirectional processing gives context from both past and future within the window.
    \item \textbf{Why both?} CNN handles local structure; BiLSTM handles temporal evolution. Together they capture patterns at multiple timescales.
\end{itemize}

\begin{lstlisting}[caption={CNN-BiLSTM price encoder (src/models/price\_model.py)}]
class PriceDirectionModel(nn.Module):
    def __init__(self, num_features, conv_channels=64,
                 lstm_hidden=128):
        super().__init__()
        # Local pattern extraction: two Conv1d layers
        self.conv = nn.Sequential(
            nn.Conv1d(num_features, conv_channels,
                      kernel_size=3, padding=1),
            nn.ReLU(),
            nn.BatchNorm1d(conv_channels),
            nn.Conv1d(conv_channels, conv_channels,
                      kernel_size=3, padding=1),
            nn.ReLU(),
            nn.BatchNorm1d(conv_channels),
        )
        # Temporal evolution: 2-layer BiLSTM
        self.lstm = nn.LSTM(
            input_size=conv_channels,
            hidden_size=lstm_hidden,
            num_layers=2,
            batch_first=True,
            dropout=0.2,
            bidirectional=True,  # 128*2 = 256-d output
        )
        # Classification head
        self.head = nn.Sequential(
            nn.Linear(lstm_hidden * 2, 128),
            nn.ReLU(),
            nn.Dropout(0.2),
            nn.Linear(128, 2),
        )

    def forward(self, x):
        # x: [B, 60, 21]
        x = x.transpose(1, 2)   # [B, 21, 60] for Conv1d
        x = self.conv(x)        # [B, 64, 60]
        x = x.transpose(1, 2)   # [B, 60, 64] for LSTM
        x, _ = self.lstm(x)     # [B, 60, 256]
        x = x[:, -1, :]         # [B, 256] last timestep
        return self.head(x)     # [B, 2]
\end{lstlisting}

\textbf{Output:} 256-dimensional embedding per stock per day.

\subsection{Model 2: Graph Attention Network (GAT)}

\subsubsection{Design Rationale}

Stocks do not move in isolation. NVIDIA's supply relationships with Apple and Microsoft, or Google's competition with Meta in digital advertising, create cross-stock dependencies that a stock-independent model cannot capture. The GAT propagates information across the 30-stock graph, allowing each stock's representation to incorporate signals from its business neighbors.

\textbf{Key constraint:} Implemented in pure PyTorch without PyTorch Geometric, as required by the project specification.

\subsubsection{GAT Attention Mechanism}

Following Veli\v{c}kovi\'c et al.\ (2018), each GAT layer computes:

\begin{align}
    e_{ij} &= \text{LeakyReLU}\left(\mathbf{a}^T [\mathbf{W}h_i \| \mathbf{W}h_j]\right) \\
    \alpha_{ij} &= \frac{\exp(e_{ij})}{\sum_{k \in \mathcal{N}(i)} \exp(e_{ik})} \\
    h'_i &= \sigma\left(\sum_{j \in \mathcal{N}(i)} \alpha_{ij} \mathbf{W} h_j\right)
\end{align}

\begin{lstlisting}[caption={Custom GAT layer (src/models/gat\_model.py)}]
class GATLayer(nn.Module):
    """Single-head graph attention layer."""
    def __init__(self, in_features, out_features,
                 negative_slope=0.2, dropout=0.1):
        super().__init__()
        self.W = nn.Linear(in_features, out_features,
                           bias=False)
        self.a = nn.Parameter(
            torch.empty(2 * out_features))
        nn.init.xavier_uniform_(self.a.unsqueeze(0))
        self.leaky_relu = nn.LeakyReLU(negative_slope)
        self.dropout = nn.Dropout(dropout)

    def forward(self, x, edge_index, edge_weight=None):
        # x: [N, in_features]
        # edge_index: [2, E] directed edges
        N = x.size(0)
        Wh = self.W(x)                  # [N, out_feat]
        src, tgt = edge_index            # each [E]

        # Attention coefficients
        Wh_src = Wh[src]                 # [E, out_feat]
        Wh_tgt = Wh[tgt]                # [E, out_feat]
        cat = torch.cat([Wh_src, Wh_tgt], dim=1)
        e = self.leaky_relu(cat @ self.a)  # [E]

        # Sparse softmax per target node
        e_max = torch.full((N,), float("-inf"),
                           device=x.device)
        e_max.scatter_reduce_(0, tgt, e,
            reduce="amax", include_self=False)
        e_exp = torch.exp(e - e_max[tgt])
        e_sum = torch.zeros(N, device=x.device)
        e_sum.scatter_add_(0, tgt, e_exp)
        alpha = e_exp / (e_sum[tgt] + 1e-10)
        alpha = self.dropout(alpha)

        # Weighted aggregation
        messages = alpha.unsqueeze(1) * Wh_src
        out = torch.zeros(N, Wh.size(1),
                          device=x.device)
        out.scatter_add_(0,
            tgt.unsqueeze(1).expand_as(messages),
            messages)
        return out
\end{lstlisting}

\begin{lstlisting}[caption={Graph-enhanced model with residual GAT (src/models/gat\_model.py)}]
class GraphEnhancedModel(nn.Module):
    def __init__(self, num_features=10,
                 encoder_dim=256, gat_hidden=64,
                 gat_heads=4, num_classes=2):
        super().__init__()
        # Shared CNN-BiLSTM price encoder
        self.conv = nn.Sequential(
            nn.Conv1d(num_features, 64, 3, padding=1),
            nn.ReLU(), nn.BatchNorm1d(64),
            nn.Conv1d(64, 64, 3, padding=1),
            nn.ReLU(), nn.BatchNorm1d(64))
        self.lstm = nn.LSTM(64, 128, num_layers=2,
            batch_first=True, dropout=0.2,
            bidirectional=True)
        # Projection + two-layer GAT with residuals
        gat_in_dim = gat_hidden * gat_heads
        self.node_proj = nn.Linear(encoder_dim,
                                    gat_in_dim)
        self.gat1 = MultiHeadGAT(gat_in_dim,
            gat_hidden, gat_heads)
        self.gat_norm1 = nn.LayerNorm(gat_in_dim)
        self.gat2 = MultiHeadGAT(gat_in_dim,
            gat_hidden, gat_heads)
        self.gat_norm2 = nn.LayerNorm(gat_in_dim)
        self.head = nn.Sequential(
            nn.Linear(gat_in_dim, 128), nn.ReLU(),
            nn.Dropout(0.2), nn.Linear(128, 2))

    def forward(self, x, edge_index, edge_weight=None):
        # x: [N, T, F] for N companies
        h = self.encode_price(x)        # [N, 256]
        h = self.node_proj(h)           # [N, gat_dim]
        # GAT layer 1 + residual + LayerNorm
        h1 = F.elu(self.gat1(h, edge_index))
        h = self.gat_norm1(h + h1)
        # GAT layer 2 + residual + LayerNorm
        h2 = F.elu(self.gat2(h, edge_index))
        h = self.gat_norm2(h + h2)
        return self.head(h)             # [N, 2]
\end{lstlisting}

\textbf{Result:} GAT contributes $+2.6\%$ AUC over a no-GAT ablation, confirming that inter-company information propagation adds discriminative signal.

\subsection{Model 3: FinBERT Document Encoder}

\begin{lstlisting}[caption={FinBERT with attention pooling (src/models/document\_model.py)}]
class AttentionPooling(nn.Module):
    def __init__(self, hidden_size):
        super().__init__()
        self.score = nn.Linear(hidden_size, 1)

    def forward(self, sequence_output, attention_mask):
        scores = self.score(sequence_output).squeeze(-1)
        scores = scores.masked_fill(
            attention_mask == 0, float("-inf"))
        weights = torch.softmax(scores, dim=-1) \
            .unsqueeze(-1)
        return torch.sum(
            sequence_output * weights, dim=1)

class DocumentDirectionModel(nn.Module):
    def __init__(self, backbone="ProsusAI/finbert",
                 dropout=0.2):
        super().__init__()
        self.encoder = AutoModel.from_pretrained(
            backbone)
        hidden = self.encoder.config.hidden_size  # 768
        self.pool = AttentionPooling(hidden)
        self.head = nn.Sequential(
            nn.Dropout(dropout),
            nn.Linear(hidden, 256),
            nn.ReLU(),
            nn.Dropout(dropout),
            nn.Linear(256, 2),
        )

    def forward(self, input_ids, attention_mask):
        output = self.encoder(
            input_ids=input_ids,
            attention_mask=attention_mask)
        pooled = self.pool(
            output.last_hidden_state, attention_mask)
        return self.head(pooled)
\end{lstlisting}

\textbf{Why frozen backbone?} With only 227 training filings, fine-tuning FinBERT's 110M parameters would cause catastrophic overfitting. Only the attention pooling layer and head (small number of parameters) are trained.

\textbf{Output:} 768-dimensional embedding per (stock, year) pair.

\subsection{Model 4: Macro MLP Encoder}

\begin{lstlisting}[caption={Macro feature encoder (src/models/macro\_model.py)}]
class MacroStateModel(nn.Module):
    """3-layer MLP: 12 -> 64 -> 64 -> 32."""
    def __init__(self, input_dim=12, hidden_dim=64,
                 output_dim=32, dropout=0.2):
        super().__init__()
        self.encoder = nn.Sequential(
            nn.Linear(input_dim, hidden_dim),
            nn.LayerNorm(hidden_dim),
            nn.GELU(),
            nn.Dropout(dropout),
            nn.Linear(hidden_dim, hidden_dim),
            nn.LayerNorm(hidden_dim),
            nn.GELU(),
            nn.Dropout(dropout),
            nn.Linear(hidden_dim, output_dim),
        )

    def encode(self, x):
        return self.encoder(x)  # [B, 12] -> [B, 32]
\end{lstlisting}

\textbf{Output:} 32-dimensional embedding encoding the current macroeconomic regime.

\subsection{Fusion Model with Gated Attention}

\subsubsection{Architecture}

The fusion model combines all four modality embeddings through a learned sigmoid gating mechanism:

\begin{equation}
    g_i = \frac{\sigma(f_{\text{gate}}(\mathbf{z}))_i \cdot m_i}{\sum_{j=1}^{4} \sigma(f_{\text{gate}}(\mathbf{z}))_j \cdot m_j + \epsilon}
\end{equation}

where $m_i \in \{0,1\}$ is the availability mask, $\sigma$ is sigmoid, and $f_{\text{gate}}$ is a two-layer MLP.

\begin{lstlisting}[caption={Phase 14 Fusion Model with HAR-RV skip (src/models/fusion\_model.py)}]
class Phase14FusionModel(nn.Module):
    """Adds HAR-RV skip connection to gated fusion."""
    NUM_MODALITIES = 4

    def __init__(self, price_dim=256, har_rv_dim=3,
                 har_proj_dim=32, gat_dim=256,
                 doc_dim=768, macro_dim=32,
                 surprise_dim=5, proj_dim=128,
                 hidden_dim=256, dropout=0.3):
        super().__init__()
        # HAR-RV skip (direct, ungated)
        self.har_rv_skip = nn.Sequential(
            nn.Linear(har_rv_dim, har_proj_dim),
            nn.LayerNorm(har_proj_dim), nn.GELU())
        # Per-modality projectors
        def _proj(in_dim):
            return nn.Sequential(
                nn.Linear(in_dim, proj_dim),
                nn.LayerNorm(proj_dim),
                nn.GELU(),
                nn.Dropout(dropout * 0.5))
        self.price_proj = _proj(price_dim)
        self.gat_proj   = _proj(gat_dim)
        self.doc_proj   = _proj(doc_dim)
        self.macro_proj = _proj(macro_dim)
        # Gating (surprise-conditioned)
        gate_in = proj_dim * 4 + surprise_dim
        self.gate = nn.Sequential(
            nn.Linear(gate_in, hidden_dim), nn.GELU(),
            nn.Linear(hidden_dim, 4), nn.Sigmoid())
        # Shared trunk: gated + HAR-RV
        trunk_in = proj_dim * 4 + har_proj_dim
        self.trunk = nn.Sequential(
            nn.Linear(trunk_in, hidden_dim),
            nn.LayerNorm(hidden_dim), nn.GELU(),
            nn.Dropout(dropout),
            nn.Linear(hidden_dim, hidden_dim),
            nn.LayerNorm(hidden_dim), nn.GELU(),
            nn.Dropout(dropout * 0.5))
        # Volatility head (Softplus ensures > 0)
        self.volatility_head = nn.Sequential(
            nn.Linear(hidden_dim, hidden_dim // 2),
            nn.GELU(), nn.Dropout(dropout * 0.5),
            nn.Linear(hidden_dim//2, hidden_dim//4),
            nn.GELU(),
            nn.Linear(hidden_dim // 4, 1),
            nn.Softplus())
        # Direction head
        self.direction_head = nn.Sequential(
            nn.Linear(hidden_dim, hidden_dim // 4),
            nn.GELU(), nn.Linear(hidden_dim//4, 2))
\end{lstlisting}

\begin{lstlisting}[caption={Forward pass with gating and HAR-RV skip}]
    def forward(self, price_emb, har_rv_raw, gat_emb,
                doc_emb, macro_emb, surprise_feat,
                modality_mask):
        p = self.price_proj(price_emb)
        g = self.gat_proj(gat_emb)
        d = self.doc_proj(doc_emb)
        m = self.macro_proj(macro_emb)
        har = self.har_rv_skip(har_rv_raw)  # [B, 32]

        stacked = torch.stack([p, g, d, m], dim=1)
        stacked = stacked * modality_mask.unsqueeze(-1)
        flat = stacked.view(stacked.size(0), -1)
        gate_in = torch.cat([flat, surprise_feat], 1)
        gates = self.gate(gate_in)
        gates = gates * modality_mask
        gates = gates / (gates.sum(1, True) + 1e-8)

        weighted = stacked * gates.unsqueeze(-1)
        fused = weighted.view(weighted.size(0), -1)
        # Concatenate gated + HAR-RV skip
        trunk_in = torch.cat([fused, har], dim=-1)
        h = self.trunk(trunk_in)

        vol = self.volatility_head(h).squeeze(-1)
        dir_logits = self.direction_head(h)
        return {"volatility_pred": vol,
                "direction_logits": dir_logits,
                "gate_weights": gates}
\end{lstlisting}

\textbf{HAR-RV Skip Connection (Phase 14 key innovation):}
The CNN-BiLSTM compresses $[60, 21] \to 256$ dimensions, losing some autoregressive signal. The skip connection extracts \texttt{rv\_lag1d}, \texttt{rv\_lag5d}, \texttt{rv\_lag22d} from the last timestep and feeds them \textit{directly} into the fusion trunk, bypassing the bottleneck. This single addition improved volatility $R^2$ from 0.867 to 0.921.

\subsection{Loss Functions}

\subsubsection{QLIKE Loss for Volatility (Weight: 0.85)}

\begin{equation}
    \mathcal{L}_{\text{QLIKE}} = \frac{1}{N}\sum_{i=1}^N \left(\frac{\sigma^2_{\text{true},i}}{\hat{\sigma}^2_i} + \log\hat{\sigma}^2_i\right)
\end{equation}

\textbf{Why QLIKE over MSE?} When $\hat{\sigma} < \sigma_{\text{true}}$ (underprediction), the first term $\sigma^2_{\text{true}}/\hat{\sigma}^2$ explodes. Underpredicting risk is more dangerous than overpredicting it. QLIKE is also scale-invariant.

\begin{lstlisting}[caption={QLIKE loss (src/models/losses.py)}]
class QLIKELoss(nn.Module):
    """QLIKE (Quasi-Likelihood) loss for volatility
    forecasting (Patton, 2011)."""
    def __init__(self, eps=1e-6):
        super().__init__()
        self.eps = eps

    def forward(self, pred, target):
        valid = ((target > self.eps)
                 & (pred > self.eps)
                 & ~torch.isnan(target))
        if not valid.any():
            return torch.tensor(0.0, device=pred.device,
                                requires_grad=True)
        pred_var = pred[valid].pow(2).clamp(min=self.eps)
        true_var = target[valid].pow(2).clamp(
            min=self.eps)
        qlike = true_var / pred_var + torch.log(pred_var)
        return qlike.mean()
\end{lstlisting}

\subsubsection{ListNet Ranking Loss for Direction (Weight: 0.15)}

\begin{equation}
    \mathcal{L}_{\text{ListNet}} = \frac{1}{|D|}\sum_{d \in D} \text{CE}\!\left(\text{softmax}(\mathbf{y}_d \cdot \tau) \;\big\|\; \text{softmax}(\hat{\mathbf{y}}_d)\right)
\end{equation}

\textbf{Why ListNet over Cross-Entropy?} Cross-entropy evaluates each stock independently. ListNet groups all 30 stocks on the same trading date and optimizes their \textit{relative ranking}, canceling out market-wide movement (beta) and forcing the model to learn stock-specific signals.

\begin{lstlisting}[caption={ListNet ranking loss (src/models/losses.py)}]
def listnet_loss(scores, labels, dates,
                 temperature=5.0):
    """ListNet ranking loss.
    Groups stocks by trading date and computes
    cross-entropy between predicted and true rankings.
    """
    pred_probs = F.softmax(scores, dim=1)[:, 1]
    loss = torch.tensor(0.0, device=scores.device)
    n_dates = 0
    for date in dates.unique():
        mask = dates == date
        n_stocks = mask.sum()
        if n_stocks < 2:
            continue
        pred = pred_probs[mask]
        true = labels[mask].float()
        pred_dist = F.softmax(pred, dim=0)
        true_dist = F.softmax(true * temperature,
                              dim=0)
        ce = -(true_dist
               * torch.log(pred_dist + 1e-8)).sum()
        loss = loss + ce
        n_dates += 1
    return loss / max(n_dates, 1)
\end{lstlisting}

\subsubsection{Combined Loss}

\begin{equation}
    \mathcal{L} = 0.85 \times \mathcal{L}_{\text{QLIKE}} + 0.15 \times \mathcal{L}_{\text{ListNet}}
\end{equation}

\textbf{Why 0.85/0.15?} Three variants were tested: balanced (0.5/0.5), dir-assisted (0.7/0.3), and vol-primary (0.85/0.15). The balanced variant stopped training at epoch 10 because QLIKE and ListNet gradients pointed in opposing directions. The 0.85/0.15 split gave one objective a dominant voice, allowing the optimizer to converge (training continued to epoch 26).

\subsection{Training Configuration}

\begin{table}[H]
\centering
\caption{Hyperparameters for the fusion model}
\begin{tabular}{lll}
\toprule
\textbf{Hyperparameter} & \textbf{Value} & \textbf{Justification} \\
\midrule
Optimizer         & AdamW            & Correct L2 regularization (vs Adam) \\
Learning rate     & $5 \times 10^{-4}$ & Standard for Adam-family on financial data \\
Weight decay      & $10^{-3}$        & L2 regularization to prevent overfitting \\
Batch size        & 2048             & Large batch for stable ListNet gradient \\
LR Scheduler      & CosineAnnealing  & Smooth decay, avoids sharp LR drops \\
Max epochs        & 60               & With early stopping \\
Early stopping    & Patience = 7     & Stops if val loss flat for 7 epochs \\
Gradient clipping & max\_norm = 1.0  & Prevents gradient explosion \\
AMP               & fp16             & Halves memory, enables larger batches \\
\bottomrule
\end{tabular}
\end{table}

\begin{lstlisting}[caption={Training infrastructure (src/train/common.py)}]
@dataclass
class TrainingConfig:
    learning_rate: float = 1e-4
    weight_decay: float = 1e-4
    epochs: int = 20
    batch_size: int = 512
    patience: int = 5
    use_amp: bool = True
    gradient_accumulation_steps: int = 1

def create_optimizer(model, config):
    return torch.optim.AdamW(
        model.parameters(),
        lr=config.learning_rate,
        weight_decay=config.weight_decay)

class EarlyStopping:
    """Stop when monitored metric stops improving."""
    def __init__(self, patience=5, mode="min",
                 min_delta=1e-4):
        self.patience = patience
        self.mode = mode
        self.min_delta = min_delta
        self.counter = 0
        self.best_score = None
        self.should_stop = False

    def __call__(self, score):
        if self.best_score is None:
            self.best_score = score
            return False
        improved = (score < self.best_score
                    - self.min_delta
                    if self.mode == "min"
                    else score > self.best_score
                    + self.min_delta)
        if improved:
            self.best_score = score
            self.counter = 0
        else:
            self.counter += 1
            if self.counter >= self.patience:
                self.should_stop = True
        return self.should_stop
\end{lstlisting}

\subsection{Results}

\subsubsection{Volatility Forecasting Performance}

\begin{table}[H]
\centering
\caption{Volatility R\textsuperscript{2} progression across all project phases}
\begin{tabular}{lrr}
\toprule
\textbf{Phase / Model} & \textbf{Vol R\textsuperscript{2}} & \textbf{Dir AUC} \\
\midrule
Historical Average (naive baseline) & 0.348 & --- \\
V1 Fusion (Phase 6)                 & 0.415 & 0.516 \\
V2 + ListNet (Phase 10)             & 0.335 & 0.599 \\
QLIKE + 30 Stocks (Phase 12)        & 0.772 & 0.568 \\
HAR-RV Features (Phase 13)          & 0.867 & 0.585 \\
\textbf{HAR-RV Skip (Phase 14)}     & \textbf{0.921} & \textbf{0.591} \\
HAR-RV Benchmark                    & 0.947 & --- \\
\bottomrule
\end{tabular}
\end{table}

\begin{figure}[H]
\centering
\includegraphics[width=\textwidth]{figures/r2_progression.png}
\caption{Volatility R\textsuperscript{2} improvement across phases. Phase 14 closes 95.7\% of the gap between V2 baseline and the HAR-RV benchmark.}
\end{figure}

\subsubsection{Learned Gate Weights}

\begin{figure}[H]
\centering
\includegraphics[width=0.7\textwidth]{figures/gate_weights.png}
\caption{Fusion gate weights in the final model. GAT (41.2\%) is dominant, confirming that cross-stock relationships are the most informative signal once data leakage is removed.}
\end{figure}

\subsubsection{Backtest Results (Direction Strategy)}

A weekly long-short strategy using the direction head:
\begin{itemize}[noitemsep]
    \item \textbf{Setup:} Top-2 long / Bottom-2 short by predicted UP probability, weekly rebalancing
    \item \textbf{Sharpe at 0bp:} 0.697 \quad \textbf{At 5bp:} 0.620 \quad \textbf{At 10bp:} 0.543
    \item \textbf{Break-even:} $\approx$20bp --- strategy remains viable at realistic transaction costs
\end{itemize}

%══════════════════════════════════════════════════════════════════════════════
\section{Summary \& Next Steps}
%══════════════════════════════════════════════════════════════════════════════

\subsection{Phase 3 Deliverables Completed}
\begin{itemize}[noitemsep]
    \item[$\checkmark$] Chronological train/validation/test split (65.8/10.9/23.3\%)
    \item[$\checkmark$] Z-score normalization (fit on training set only)
    \item[$\checkmark$] 21-feature engineering pipeline including HAR-RV autoregressive features
    \item[$\checkmark$] FinBERT chunking pipeline for 227 SEC filings
    \item[$\checkmark$] Macroeconomic feature extraction with 1-day lag
    \item[$\checkmark$] 30-node knowledge graph construction
    \item[$\checkmark$] Availability masking for missing modalities
    \item[$\checkmark$] Two data leakage audits and corrections
\end{itemize}

\subsection{Phase 4 Deliverables Completed}
\begin{itemize}[noitemsep]
    \item[$\checkmark$] CNN-BiLSTM price encoder (256-d output)
    \item[$\checkmark$] Custom GAT in pure PyTorch (256-d output, +2.6\% AUC over ablation)
    \item[$\checkmark$] Frozen FinBERT with attention pooling (768-d output)
    \item[$\checkmark$] Macro MLP encoder (32-d output)
    \item[$\checkmark$] Gated multimodal fusion with availability masking
    \item[$\checkmark$] HAR-RV skip connection ($R^2$: 0.867 $\to$ 0.921)
    \item[$\checkmark$] QLIKE loss for volatility + ListNet ranking loss for direction
    \item[$\checkmark$] AMP fp16 training with gradient clipping
\end{itemize}

\subsection{Key Lessons Learned}
\begin{enumerate}
    \item \textbf{Data integrity is paramount.} Two separate leakages demonstrated how easily an ML pipeline can appear to work when a feature encodes the target. Auditing inputs when results seem too good is essential.
    \item \textbf{Task choice matters.} Volatility is more predictable than direction on efficient markets. Switching the primary objective from direction to volatility increased $R^2$ from 0.335 to 0.921.
    \item \textbf{Skip connections preserve information.} The HAR-RV skip connection directly addressed the CNN-BiLSTM compression bottleneck, providing the single largest improvement ($+6.3\%$ $R^2$).
    \item \textbf{Loss function design is architecture.} QLIKE's asymmetric penalty and scale-invariance, and ListNet's beta-canceling ranking formulation, contributed as much as architectural changes.
\end{enumerate}

\vspace{1cm}
\noindent\textbf{GitHub Repository:} \url{https://github.com/imrishi007/financial-document-analysis}
